\begin{figure}[htbp]
\begin{center}
\includegraphics[width=0.9\textwidth]{FlxPin_preopt_20261026.eps}
\end{center}
\caption{Pre-optimization pinned-knee torque versus knee angle for the four cross-validation flexor tests. Knee angle is negative in flexion, and flexion torque is negative. Gold circles are the hybrid calculation, which combines the measured muscle length and moment arm with the modeled force; indigo filled circles are experimentally measured torque; the gray dashed line is the Hunt-equation rigid-body prediction, evaluated directly from the equation of record and defined only where its modeled strain is non-negative; and the light orange dash-dot line is the unoptimized prediction, the improved BPA force characterization applied to the same rigid kinematics. Panel titles give the nominal resting length; the actual resting lengths are $l_{0}=\qty{48.5}{\cm}$ \textbf{(A)}, $l_{0}=\qty{41.0}{\cm}$ \textbf{(B)}, $l_{0}=\qty{45.7}{\cm}$ \textbf{(C)}, and $l_{0}=\qty{40.6}{\cm}$ with tendon \textbf{(D)}. Panels \textbf{(A)} and \textbf{(B)} are the training tests and panels \textbf{(C)} and \textbf{(D)} the validation tests of the adopted cross-validation fold.}\label{fig:FlxPin_preopt}
\end{figure}
